% 原书第 65 页（PDF 第 78 页）
\BookPage{65}{78}
\BookHeading{扬州二十四桥考}{扬州二十四桥考}{何庄*}
唐文宗大和七年至九年（833-835）著名诗人杜牧任扬州节度使掌书记，与同僚韩绰交好，离任后不久，他写下了一首赠寄之作〈寄扬州韩绰判官〉。全文如下:「青山隐隐水迢迢，秋尽江南草未凋。二十四桥明月夜，玉人何处教吹箫。」这首诗描绘了扬州动人的风貌，探问了好友的近况，意境优美，风调悠扬，是杜牧以扬州为题材的佳作中最著名的篇章。一经问世，二十四桥的美名便不胫而走，千古传诵。二十四桥成了扬州这个美丽城市的象征，成了扬州的名片，凡到扬州者，无不寻访二十四桥的踪迹。同时，这首诗也留下了一椿千年疑案：诗中的二十四桥究竟指的是什么？是二十四座桥，还是一座桥，或指其他？千百年来，无数的文人学者苦苦追寻，但是众说纷纭，至今尚无定论。综合起来，主要有以下几种说法: 1.二十四桥说。此说以北宋科学家沈括的《梦溪笔谈》和南宋祝穆的《方舆胜览》为代表；2.一桥说。此说以清初历史学家谈迁的《北游记》和乾隆时期李斗的《扬州画舫录》等为代表；3.编号说。有人认为，二十四桥是扬州城里编号为二十四的一座桥；4.谐音说。近年来，一位扬州的青年作者从扬州方言的角度考证，二十四桥即「阿师桥」的谐音；5.泛指说。不少人认为杜牧诗中的二十四桥是文学的惯用手法，泛指扬州水道纵横，桥梁众多，不必拘泥。2五种说法中，第二、三、四种说法总体上可归于一桥说，只是具体解释不同而已，实际上主要就是三种说法。

笔者翻阅了大量有关的考证文章，发现大多是从诗歌出发进行史实考证的，就是以杜牧的诗为既定史实，考证二十四桥的兴废变迁，或者将史实的考证与诗歌本意的考证混为一谈。笔者认为，要考证二十四桥，必须首先区别开杜牧诗中二十四

\vfill
\begin{BookNotes}
\BookNote{*}{作者任职于中国人民大学档案学院。}
\BookNote{1.}{朱千华，〈二十四桥今何在〉，《中华读书报》1998年8月25日。}
\BookNote{2.}{吴锡平，〈扬州二十四桥〉，《寻根》2003(5)。}
\end{BookNotes}

% 原书第 66 页（PDF 第 79 页）
\BookPage{66}{79}
\noindent 桥的本意和历史上真实的二十四桥。这是不同的概念，分属于不同的领域，前者属于文学，后者属于史学。一方面，诗歌是一种艺术，艺术源于生活，并非完全凭空想像，而是有所依据，特别是具有史诗性质的诗歌。所以诗歌可以作为史料用来证史之确、补史之阙或纠史之误，诗文证史无疑是可行的。以杜牧诗为例，杜牧〈扬州三首〉诗中所提到的雷塘、迷楼、斗鸡台，均为有据可查的古迹，诗中提到的「霞映两重城」的内城和外城，也透露出唐代扬州城分为「子城」和「罗城」的事实，据高彦林《唐阙史》载：「扬州，胜地也，每重城向夕…九里三十步街中，珠翠塡咽。」杜牧诗可以印证这一史实；同时，也可以借助历史知识加深对诗文的理解，史解诗文也是可行的，比如杜牧〈赤壁〉中提到「东风不予周郎便」，用的是三国时期周瑜火烧赤壁、成就孙权霸业的典故，可以借助《资治通鉴》和《三国志·周瑜传》中有关历史记载理解诗句，即所谓诗史互证。在这方面，近代考据学大师陈寅恪堪称典范。他的《元白诗笺证稿》精彩之处，在于既能「史解诗文」，又能「以诗证史」，把文学的研究与历史的研究很好地结合起来。

另一方面，艺术又高于生活，需要对生活进行必要的夸张或浓缩，使情感的表达更为强烈和集中，因此，诗歌中的史实是生动形象的，也是概要的，不能把诗歌中的描述当成准确的史实，而历史的考证，必须以真实准确的史料为基础，这是在运用诗歌史料时特别需要注意的。就上述问题而言，如果拿杜牧诗歌中的二十四桥一一对应现实，诗中寻史，因诗名桥，作为文人怀旧的风雅之举倒也无可厚非，但是就严谨的历史考证而言，则违反了史学的基本精神。因此，笔者以下对于二十四桥的的考证分为两部分，第一部分是杜牧诗中本意的考证，第二部分是结合各说对史实的考证，不妥之处，敬请各位专家指正。

\BookSubheading{一、二十四桥本义考}

论及杜牧诗中二十四桥的的本意，首先要考察杜牧「二十四」这个数字的用法。如同要解决《红楼梦》的作者问题，曹雪芹在书中的表白是最重要的证据一样；要知道「二十四」何指，从杜牧诗歌中寻找本证当为首要。根据岳麓书社集部丛刊《杜牧集》的统计，杜牧诗歌450余首，其中有数字的约265首，超过了半


% 原书第 67 页（PDF 第 80 页）
\BookPage{67}{80}
\noindent 数，为其他诗人所罕见。所以杜牧善于以数字入诗，在后代是著名的。如何文焕在《历代诗话》中说:「旧称牧之好用数目，如二十四桥之类是也。」杨慎在《升庵诗话》中也说:「大抵牧之诗好用数目。」翻检杜牧的诗作，以数字入诗的情况，可分为两种:一种情况是实指，即诗中数字与实际完全吻合。这种情况不太多，如〈感怀诗〉的「九朝仗神灵，四海为输委，如何七十年，汗含羞耻」，九朝实指皇帝的祖庙，七十年实指安史之乱爆发到李同捷谋反这段时间。〈自遣〉中「四十已云老，况逢忧窘余」，时年作者四十岁。〈春末题池州弄水池〉的「使君四十四，两佩左铜鱼」，时年作者四十四岁。新定途中〉的「无端偶效张文纪，下杜乡园别五秋」，五秋实指杜牧在会昌二年离长安赴黄州任刺史，正好五年。还有〈题齐安城楼〉的「不用凭栏苦回首，故乡七十五长亭」据杜佑《通典》云，齐安郡「去西京二千二百五十五里」，刚好是七十五亭。

第二种情况是泛指，即诗中数字与实际大致相当。这种情况占大多数。如〈江南春绝句〉「南朝四百八十寺，多少楼台烟雨中」，据《南史·郭祖深传》「都下佛寺，五百余所，穷极宏丽」，四百八十寺应该是个约数。〈隋堤柳〉「夹岸垂杨三百里，只应画里最相宜」，指的是两岸栽满杨柳，连绵数百里。〈村舍燕〉「汉宫一百四十五，多下珠帘闭锁窗」，引用了张衡〈西京赋〉中「封畿千里，郡国宫管馆，百四十五」的说法，据《文选》中李善注引的〈三辅故事〉云，张衡套用的是「秦时殿馆，百四十五所」的记载，可见一百四十五也是约数。〈题宣州开元寺〉「楼飞九十尺，廊环四百柱」，显然是约数，形容寺楼之高大。再如〈题敬安寺楼〉「暮景千山雪，春寒百尺楼」，不用说也是约数，描绘高楼之雄伟，雪景之壮阔。

笔者分析发现，杜牧的数字入诗有规律可循。数字实指的情况（时间、年龄、路程等）多为纪实性的诗篇，即杜牧个人经历或有关史实；泛指的情况（自然景观、人文古迹等）多为杜牧描写风物的抒情之作，而〈寄扬州韩绰判官〉这首诗无疑属于后者，所以，二十四桥理应代指扬州水道纵横，桥梁众多，不会指某一座特定的桥。另外，仅从这首诗的表达看，「玉人何处教吹箫」这句中「何处」一词，表明不止一处，当然不会是一座桥了，这就是我们从本证得出的结论。事实上，泛


% 原书第 68 页（PDF 第 81 页）
\BookPage{68}{81}
\noindent 指是诗歌最常见的现象，正体现了艺术源于生活，又高于生活的特征。不仅杜牧的诗歌如此，其他诗人吟咏风物的诗歌亦如此。如白居易描绘唐代苏州桥梁云：「江栏三百九十桥。」南宋杨备则云:「画桥四百」，而南宋《平江图》标出的苏州桥梁的准确数目是359座，诗歌中的数字同样是一个泛指的约数。

清楚了杜牧的二十四桥所指，进一步说，二十四桥的本意仅仅指的是扬州优美的夜景吗？笔者认为不是，它应该是泛指扬州的繁华，尤其是指歌弦声喧的「烟花月夜」。最充分的证明仍然在杜牧本人的诗集中。杜牧写扬州的诗，还有〈赠别两首，其一云:「娉娉嫘嫘十三余，豆蔻梢头二月初。春风十里扬州路，卷上珠帘总不如。」其二云:「多情却是总无情，唯觉尊前笑不成。蜡烛有心还惜别，替人垂泪到天明。」这两首诗都很著名，广为流传，是杜牧写给扬州歌伎的风情之作，是他被称为「风流才子」的证据，对此，杜牧本人并不否认，他在〈遣怀〉中写道:「落魄江南载酒行，楚腰纤细掌中轻。十年一觉扬州梦，赢得青楼薄幸名。」

轻狂是当时一般文人的通病，笔者并不想为杜牧掩饰。需要说明的是，杜牧对扬州的向往和留恋，更多的不是那里的美景，而是那里的美人。况且，自古就有扬州出美女的说法，《红楼梦》里的林黛玉，不就是扬州的吗？如果把「二十四桥明月夜」和「玉人何处教吹箫」联系起来，画桥丽影，这层意思就更清晰了，意境也更优美了。中华民族的审美习惯，向来是马上看英雄，月下看美人，杜牧的诗句可以看作是这一审美习惯的注脚。历代文人对二十四桥的描写，无不着重于「烟花月夜」这一层意思，追寻二十四桥，就是追寻繁华之梦。我们不妨略举几例。南宋著名词人姜夔的〈扬州慢·淮左名都〉:「二十四桥仍在，波心荡，冷月无声。念桥边红药，年年知为谁生!」联系前面的几句:「杜郎俊赏，算而今，重到须惊。纵豆蔻词工，青楼梦好，难赋深情。」诗中意境，显然是借桥念及桥边的玉人，玉人不在，红药何用。与姜夔同时代的另外几位诗人，如韩琦：「二十四桥千步柳，春风十里上珠帘。」吴文英：「二十四桥南北，罗存香分。」清代的诗人也不例外，杨夔生〈扬州好〉：「扬州好，二十四桥头，花扑酒帘风扬恨，月随歌舫水添愁。等等。对于二十四桥的本意，明末清初著名文学家、史学家张岱的理解应该是最准确的。他在〈陶庵梦忆·二十四桥风月〉中写道：


% 原书第 69 页（PDF 第 82 页）
\BookPage{69}{82}
\begin{BookQuote}
广陵二十四桥风月，邗沟尚存其意。渡钞关，横互半里许，为巷者九条。巷故九，凡周旋折旋于巷之左右前后者，什百之。巷口狭而肠曲，寸寸节节，有精房密户，名妓、歪妓杂处之。名妓匿不见人，非向导莫得入。歪妓多可五六百人，每日傍晚，膏沐熏烧，出巷口，倚徙盘礴于茶馆酒肆之前，谓之「站关」。茶馆酒肆岸上纱灯百盏，诸妓掩映闪灭于其间，疤戾者帘，雄趾者阈。灯前月下，人无正色，所谓「一白能遮百丑」者，粉之力也。游子过客，往来如梭，摩睛相觑，有当意者，逼前牵之去；而是妓忽出身分，肃客先行，自缓
\end{BookQuote}

步尾之。

可见，二十四桥其实是一个典型的风月场所，类似于南京的秦淮河。我们还可以从地名和地理位置上找到二十四桥与风月联系的证据。从春秋末期吴王夫差在蜀冈建邗城起，直到西元八、九世纪的唐代，扬州在蜀冈的位置基本没变，只是在原有城市的基础上，在蜀冈下开始筑罗城，形成工商业者聚居的「十里长街」。经历了明末清初的战乱后，扬州在清乾隆时期社会经济达到鼎盛，重新成为繁华的江南胜地，写于这一时期李斗的《扬州画舫录》详尽地描绘了扬州的历史风物，展现了一个人间天堂般美丽而富庶的扬州。他的记载从地理上是顺序的，第九卷恰恰有个「小秦淮」，其名来自南京的秦淮河，多记街市之繁华、歌伎之故事，接着几卷便是虹桥及虹桥东、虹桥西、蜀冈东、蜀冈西，二十四桥在第十五卷蜀冈西。可见二十四桥正好在「小秦淮」这一带，也就是杜牧诗歌中的「春风十里扬州路」。

\BookSubheading{二、二十四桥史实考}

关于扬州二十四桥的史实是怎样的？目前依然没有定说，以下对各种说法加以辨析。辨析的目的是搞清哪种说法的可信度最高。最早记载扬州二十四桥是在隋代。隋炀帝开凿大运河，使处于水运枢纽的扬州日益重要，扬州的桥梁开始遍布全城。据清《一统志》载:「扬州二十四桥，在府城，隋置，并以城门坊市为名。」志书兼记各地地理与历史，是中国传统的书籍门类，具有很高的史料价值，清代的《一统志》为最丰富、最完善的全国性总志，不仅记载清代的历史地理，而且兼及前代，是研究中国历史地理的必读书目，其可信度毋庸置疑。二十四桥既以城门坊


% 原书第 70 页（PDF 第 83 页）
\BookPage{70}{83}
\noindent 市为名，自然不可能是一座桥了。唐代关于扬州的史料少而简，仅有上述《唐阙史》等提及扬州城的大致规模，不及二十四桥。杜牧写入诗中是二十四桥，估计唐代保持隋置未变。北宋科学家沈括是第一个详尽考证二十四桥的人。他在《梦溪笔谈·补笔谈》卷三中云:「可纪者有二十四桥。最西浊河茶园桥、次东大明桥、入西水门有九曲桥，次东正当帅牙南门有下马桥，又东作坊桥，桥东河转向南，有洗马桥，次南桥，又南阿师桥…。」将二十四桥的名称一一列出，并注明当时存在的桥仅有八座了。沈括是著名的科学家，治学严谨，他的考证应该是可信度很高的。

南宋祝穆的志书《方舆胜览》为宋代名著佳志，依次记载了南宋时期各地的地理与历史，在其中的〈淮东路·扬州〉中记有二十四桥云：「隋置，并以城门坊市为名，后韩令坤省筑州城，分布阡陌，别立桥梁，所谓二十四桥者，或存或废，不可得而考。」下面还附有杜牧、苏轼和欧阳修描写二十四桥的诗。书中沿袭了《一统志》关于二十四桥隋置的说法，「所谓二十四桥者，或存或废，不可得而考」。说明到了南宋，二十四桥的名称、数目已经有所变化，而且无从考证了，这与沈括的结果基本一致，也是比较可信的说法，但缺乏旁证。

由于元代的《一统志》在元末战乱中散佚，仅于《永乐大典》、《明一统志》、《满洲源流考》等诸书散见其佚文，故元代对于扬州城的记载很少，目前仅有《马可波罗纪行》记载了扬州社会的种种形貌，书中说道:「从泰州发足，向东南骑行一日，终抵扬州。城甚广大，所属二十七城，皆良城也。…。」但于扬州的桥梁等未及。明代嘉靖年间盛仪编纂的《维扬志》38卷，体例完善，详载扬州沿革、寺院、河流、城门、琼花、人物等。尤其是宋代编撰的扬州三志散佚，这部关于扬州城的专志成为考证扬州历史的珍贵资料，惜于津梁部分缺载。明代还有一部齐东野人编纂的《隋炀帝艳史》，记载隋炀帝与萧后赏月的故事。说炀帝与萧后来到一座新桥，见天空月光皎洁，桥上箫声悠扬。桥名末定，萧后请炀帝赐名。因同游者有二十四位宫女，遂命名为二十四桥，齐东野人为小说家，关于二十四桥的描绘想像丰富，不足为信。所以尚未发现元明时期关于扬州二十四桥的可信记载。

清初历史学家谈迁治学极为严谨，为完成明代编年史《国榷》，他多方搜集史料，考订史实，《北游录》即他北上搜集史料时留下的珍贵日记。日记详载他所


% 原书第 71 页（PDF 第 84 页）
\BookPage{71}{84}
\noindent 经之地的所见、所闻，其中有他寻访二十四桥的经过。顺治十年七月十二日（乙已），谈迁游览了扬州西郊的梅花岭、八角亭、天甯寺、司徒庙，「出问二十四桥，云南且里许」，因此作诗一首〈二十四桥〉以记其游：「斜阳古道接轮蹄，明月扶疏万柳西。桥上行人桥下水，落花尚自怨香泥。」从谈迁的游记和诗歌中看，二十四桥应该是一座桥的形象，谈迁的记载是相当可信的。

清代吴绮在《扬州鼓吹词·序》中说：「出西郭二里许，有小桥，朱栏碧，题曰烟花夜月，相传为二十四桥旧址，盖本一桥，会集二十四桥美人于此，故名。」由于它的记载基于「相传」，所以可信度不高。李斗在他的《扬州画舫录》中即认为吴绮的「美人之说，盖附会言之矣」。他这样记载：「二十四桥，即吴家砖桥，一名红药桥，在春熙台后，平泉涌瀑之水，即金匮山水，由二十四桥而来也，桥跨西门街东西两岸，砖墙庋版，围以红栏，直西通新教场，北折入金匮山，桥西吴家瓦屋圩墙上，石刻烟花夜月四字，不著书者姓名。」记载乾隆时期扬州社会、经济、文化的书籍很多，李斗的《扬州画舫录》是其中的佼佼者。它属于笔记体，仿《水经注》体例，按地分卷，巧妙地将与之相关的历史、文化、掌故、风情穿插其间，兼具史料价值与可读性，为研究扬州历史，尤其清代扬州的必读书目，他的记载也是可信的。民国时期王振世的《扬州揽胜录》基本沿袭了李斗的说法，此不赘述。

嘉庆年间由阿克当阿修、姚文田等编纂的《重修扬州府志》是一部扬州的名志，甚为学界所推崇，《续修四库全书》称「考证精详，征引宏富，而且笔严简，具有良史之法，洵称杰作」。其中卷十七〈津梁志〉首列二十四桥，记载较详：「二十四桥，隋置，在府旧城。《方舆胜览》云二十四桥并以城门坊市为名，后韩令坤省筑州城，分布阡陌，别立桥梁，所谓二十四桥者，或存或废，不可得而考。沈括补笔谈云扬州载唐时最为富盛，旧城南北十五里一百一十步，东西七里三十步可纪者有二十四桥。…案《方舆胜览》既云存废不可得而考，补笔谈所载何历历可数也？又传炀帝于月夜同宫女二十四人吹箫桥上因名。则所谓二十四桥者止一桥矣！…旧闻互异并存以俟考古者。」可以看出，这部名志确如《续修四库全书》所言，由于时代久远和众说纷纭，只好综合各说，存疑待证，不敢妄下结论。


% 原书第 72 页（PDF 第 85 页）
\BookPage{72}{85}
综上所论，笔者认为，二十四桥的史实经历了一个从唐代经宋、元、明到清代由多变少的过程，应该是比较合理的。前文已经考证，杜牧诗中的二十四桥是一个约数，泛指桥多。如果唐代扬州果真有二十几座桥，那么经过战乱，到了北宋，或废弃，或湮没，数目减少，是可以理解的。所以沈括在有些桥名下注「今存」字样。从唐代和北宋的诗歌可以明显地发现，这一时期诗歌中的二十四桥泛指多数，如欧阳修：「都将二十四桥月，换得西湖十顷秋。」黄庭坚：「淮南二十四桥月，马上时时梦见之。」苏轼：「二十四桥亦何有，换此十顷玻璃风。」而南宋以降，多指一桥，如前述姜夔、韩琦、吴文英等。说明金兵南下，扬州城被毁，原有的二十四桥，也许只剩下一座了。明代的扬州，是群雄争霸、战火不断的时期，扬州城屡遭劫难，故程文德的诗云：「倒看飞鸟穿林外，俯听吹箫出境中，二十四桥都不见，更从何处见离宫？」可能扬州的桥梁都损坏了。清代无论官方、私人记载，还是诗歌均指一桥或综合之前各说。

以上关于二十四桥的考证为地上部分，是谓书证，而其准确的数目还有待考古发掘的物证。近年来，扬州城的考古工作不断进展，古代的城址、桥梁等均有发现，如2009年发现的宋代平山堂城遗址，3还有始建于宋代的象鼻桥等等，4特别是2008年在瘦西湖一带发掘出唐代扬州城墙的「西门遗址」，发现了四座城门，为唐代的「罗城」，直接印证了唐代商业的繁华，同时这一带也是杜牧诗中二十四桥的位置，或许将来的某一天，考古提供的最新资料会揭开二十四桥的真正谜底，让我们拭目以待。

\vfill
\begin{BookNotes}
\BookNote{3.}{《扬州时报》2009年6月11日，通讯员：薛炳宏等。}
\BookNote{4.}{同上。}
\BookNote{5.}{《扬子晚报》2008年12月10日，葛学涛、陈咏。}
\end{BookNotes}
